% CM_ws paper-alignment formula patch, 2026-08-24.
% This is a manuscript snippet, not a standalone LaTeX document.

% -----------------------------------------------------------------------------
% 1. Replace the first two terms of the DNMPC objective with the following.
% W_di remains the terminal Cartesian tracking term of the local reference.
% W_ti uses the communicated predicted position of the peer EE.
% -----------------------------------------------------------------------------

\begin{align}
\min_{\boldsymbol{u}_i(t)} E_i(t)
={}&
\left\|
\boldsymbol{x}_{\mathrm{p}i}(t_0+t_{\mathrm d})
-\boldsymbol{x}^{*}_{\mathrm{p}i}(t_0+t_{\mathrm d})
\right\|_{W_{\mathrm d i}(t)}^2
\nonumber\\
&+\int_{t_0}^{t_0+t_{\mathrm d}}
\Big(
\left\|
\boldsymbol{x}_{\mathrm{p}i}(t)
-\boldsymbol{x}'_{\mathrm{p}i}(t)
\right\|_{W_{\mathrm t i}(t)}^2
\nonumber\\
&\qquad
+\left\|
\boldsymbol{x}_{\mathrm{o}i}(t)
+\boldsymbol{x}'_{\mathrm{o}i}(t)
\right\|_{W_{\mathrm o i}(t)}^2
+\left\|\boldsymbol{u}_i(t)\right\|_{W_{\mathrm u i}(t)}^2
\nonumber\\
&\qquad
+\left\|
\boldsymbol{q}_i(t)-\boldsymbol{P}^{*}_i(t)
\right\|_{W_{\mathrm q i}(t)}^2
+\left\|
\dot{\boldsymbol{q}}_i(t)-\dot{\boldsymbol{P}}^{*}_i(t)
\right\|_{W_{\mathrm v i}(t)}^2
\nonumber\\
&\qquad
+\left\|
\boldsymbol{x}_i(t)
-\boldsymbol{f}_i\bigl(\boldsymbol{x}_i(t),
\boldsymbol{x}'_i(t)\bigr)
\right\|_{W_{\mathrm F i}(t)}^2
\Big)\,\mathrm dt .
\label{eq:MPC_cost_function_revised}
\end{align}

The term weighted by $W_{\mathrm d i}(t)$ is the terminal Cartesian
tracking error with respect to the local reference trajectory. The term
weighted by $W_{\mathrm t i}(t)$ coordinates relative position by driving the
$i$th EE toward the communicated predicted position
$\boldsymbol{x}'_{\mathrm p i}(t)$ of the peer EE. Together with the joint
trajectory tracking and adaptive internal-force terms, it produces the desired
bilateral clamping tendency without replacing the generated joint reference.

% -----------------------------------------------------------------------------
% 2. Define the projected internal force used by AIC.
% Forces are first transformed to the common frame and gravity-compensated.
% -----------------------------------------------------------------------------

\begin{align}
\boldsymbol{n}_{\mathrm{lr}}(t)
&=
\frac{\boldsymbol{x}_{\mathrm{pr}}(t)
-\boldsymbol{x}_{\mathrm{pl}}(t)}
{\left\|\boldsymbol{x}_{\mathrm{pr}}(t)
-\boldsymbol{x}_{\mathrm{pl}}(t)\right\|},\\
F_{\mathrm{int},\mathrm l}(t)
&=
\left|
\boldsymbol{n}_{\mathrm{lr}}^{\top}(t)
\boldsymbol{F}_{\mathrm l}^{\mathrm c}(t)
\right|,\\
F_{\mathrm{int},\mathrm r}(t)
&=
\left|
-\boldsymbol{n}_{\mathrm{lr}}^{\top}(t)
\boldsymbol{F}_{\mathrm r}^{\mathrm c}(t)
\right|,\\
\Delta F_i(t)&=F_{\mathrm{int},i}(t)-F_{\mathrm d i}.
\label{eq:projected_internal_force}
\end{align}

Here, $\boldsymbol{F}_{i}^{\mathrm c}$ is the gravity-compensated force
expressed in the common frame. In the deformation and adaptive-control
equations, $\|\boldsymbol{F}_i\|$ is replaced by
$F_{\mathrm{int},i}$.

% -----------------------------------------------------------------------------
% 3. Clarify the force treatment over the prediction horizon.
% The measured force is held constant; geometry is predicted at every step.
% -----------------------------------------------------------------------------

For one NMPC update, the measured internal force is held constant over the
prediction horizon,
\begin{equation}
\bar F_{\mathrm{int},i}(k)=F_{\mathrm{int},i}(t_0),
\qquad k=0,\ldots,N_{\mathrm p}-1.
\end{equation}
The current adaptive variables $\hat b_i(t_0)$ and $\hat k_i(t_0)$ are also
held constant during this horizon and are updated again at the next measured
control cycle. In contrast, the deformation geometry is recomputed from the
predicted EE positions at every prediction step:
\begin{align}
d_i(k)
&=\left\|
\boldsymbol{x}'_{\mathrm p i}(k)
-\boldsymbol{x}_{\mathrm p i}(k)
\right\|,\\
e_i(k)&=\frac{1}{2}\bigl(d_{\mathrm t}-d_i(k)\bigr),\\
\dot e_i(k)&=\frac{e_i(k)-e_i(k-1)}{T_{\mathrm s}},\\
e_{i,\mathrm d}(k)
&=\operatorname{sat}_{e_{\max}}
\left(
\frac{
\bar F_{\mathrm{int},i}(k)-F_{\mathrm d i}
-\hat b_i(t_0)\dot e_i(k)}
{\hat k_i(t_0)}
\right),\\
\boldsymbol{n}_i(k)
&=\frac{
\boldsymbol{x}'_{\mathrm p i}(k)
-\boldsymbol{x}_{\mathrm p i}(k)}
{d_i(k)},\\
\boldsymbol{f}_{\mathrm p i}(k)
&=\boldsymbol{x}_{\mathrm p i}(k)
+\boldsymbol{n}_i(k)
\operatorname{sat}_{e_{\max}}
\bigl(e_{i,\mathrm d}(k)-e_i(k)\bigr).
\label{eq:aic_horizon_geometry}
\end{align}
A small nonzero denominator is used when
$|\hat k_i(t_0)|$ approaches zero. This formulation does not introduce an
unmodelled future-force predictor: it uses zero-order hold for the measured
force and the communicated position predictions for the future geometry.

% -----------------------------------------------------------------------------
% 4. Replace the signed kinematic penalties by nonnegative violations.
% These equations reproduce the current CM_ws implementation.
% -----------------------------------------------------------------------------

Define $[z]_{+}=\max(0,z)$. The joint-position penalty is
\begin{equation}
J_{\mathrm q i,m}
=\sum_{k=1}^{N_i}
\begin{cases}
P_{im,k}^{2}(t)-q_{ik,\max}^{2},
&P_{im,k}(t)>q_{ik,\max},\\
P_{im,k}^{2}(t)-q_{ik,\min}^{2},
&P_{im,k}(t)<q_{ik,\min},\\
0,&\text{otherwise},
\end{cases}
\label{eq:MINCO_joint_limit_revised}
\end{equation}
and, for symmetric velocity limits, the joint-velocity penalty is
\begin{equation}
J_{\mathrm v i,m}
=\sum_{k=1}^{N_i}
\left[
\dot P_{im,k}^{2}(t)-\dot q_{ik,\max}^{2}
\right]_{+}.
\label{eq:MINCO_velocity_limit_revised}
\end{equation}

For either grasp or transport, a relative-distance interval
$d_{\min}\leq d\leq d_{\max}$ is penalized by
\begin{equation}
J_{\mathrm p}
=\left[d^{2}-d_{\max}^{2}\right]_{+}
+\left[d_{\min}^{2}-d^{2}\right]_{+}.
\label{eq:MINCO_distance_penalty_revised}
\end{equation}
The signed orientation penalty is
\begin{equation}
J_{\mathrm o}
=\left[
\cos\theta_{\mathrm d}\,\|\boldsymbol p\|
-\boldsymbol{x}_{\mathrm o}^{\top}\boldsymbol p
\right]_{+}^{2},
\label{eq:MINCO_orientation_penalty_revised}
\end{equation}
which rejects an EE axis that points in the opposite direction.

% -----------------------------------------------------------------------------
% 5. Replace the manuscript description of J_ci and collision handling.
% -----------------------------------------------------------------------------

The components $J_{\mathrm p i}$, $J_{\mathrm u i}$,
$J_{\mathrm K i}$, $J_{\mathrm c i}$, and $J_{\mathrm T}$ denote the
trajectory-length cost, joint-velocity cost, kinematic-constraint violation
cost, task-constraint violation cost, and time cost, respectively.

Collision constraints are enforced during the preceding sampling-based path
generation. To reduce optimization time, no collision penalty, collision
gradient, or collision projection is included in the trajectory-smoothing
optimization. The optimized trajectory can be checked by FCL, but a failed
trajectory is not projected back into the collision-free set. This design
improves computational efficiency at the cost of lower collision robustness
during path smoothing.

